\documentclass{article}
\usepackage{graphicx} % Required for inserting images
\usepackage{booktabs}
\usepackage{makecell}
\usepackage{tabularx}

\title{Compte-rendu des TPs de\\Simulation d'Atomistique Avancé}
\author{Joanny Magat}
\date{September 2026}

\begin{document}

\maketitle

\section{Introduction}

\section{TP4 (02/10/2026)}

\subsection{Objecif du TP}

\subsection{Vectorisation : boucles}

\vspace{0.5cm}

\begin{tabularx}{\textwidth}{c c X X}
    \toprule
    n° de Boucle & Vectorisée ? & Commentaire & Point bloquant\\
    \midrule
    0 & Oui & optimized: loop vectorized using 16 byte vectors & x\\
    1 & Oui & optimized: loop vectorized using 16 byte vectors & x\\
    2 & Oui & optimized: loop vectorized using 8 byte vectors & x\\
    3 & ??? & rien n'est dit concernant cette boucle & x\\
    4 & Oui & optimized: loop vectorized using 8 byte vectors & x\\
    5 & Oui & optimized: loop vectorized using 16 byte vectors & x\\
    6 & Oui & optimized: loop vectorized using 16 byte vectors & x\\
    7 & Non & missed: couldn't vectorize loop & \makecell[l]{On n'a pas forcément $d > L$ \\ car l'indice est dynamique} \\
    \bottomrule
\end{tabularx}

\vspace{0.5cm}

gcc -Q --help=target \\
msse2 [enabled] \\
donc 128 bytes \\


\subsection{Vectorisation : re-indexation}

Concernant la boucle calculant la somme on a :\\
"vectorisation-reindexation.c:25:18: missed: couldn't vectorize loop"\\

\vspace{0.5cm}

Il faut donc faire une liste A_inv contenant les éléments de A trié dans l'ordre décroissant.\\
De plus, il faut transformer S en liste dont chaque élément sera à sommer.\\
La boucle deviendra alors verctorisable car tous ses éléments seront fonction de i.\\

Le code devient : \\

#include <stdio.h>
#include <assert.h>

int main() {	  

  const int dim=100000;
  double A[dim];
  double A_inv[dim];

  // initialise array
  for (int i=0; i<dim; i++) {
    A[i] = 1./(i+1);
    
  }
  for (int i=0; i<dim; i++) {
    A_inv[dim-i-1] = A[i];
  }

  // calculate required sum
  double S[dim];
  for (int i=0; i<dim; i++)
    S[i] = (i+1)*A[i] +1./(i+1)*A_inv[i]; 
  
  double sum = 0;
  for (int i=0; i<dim; i++)
    sum += S[i];
  printf("S: %lf",sum); // dummy print
}
\end{document}